\IfFileExists{stacks-project.cls}{%
\documentclass{stacks-project}
}{%
\documentclass{amsart}
}

% For dealing with references we use the comment environment
\usepackage{verbatim}
\newenvironment{reference}{\comment}{\endcomment}
%\newenvironment{reference}{}{}
\newenvironment{slogan}{\comment}{\endcomment}
\newenvironment{history}{\comment}{\endcomment}

% For commutative diagrams we use Xy-pic
\usepackage[all]{xy}

% We use 2cell for 2-commutative diagrams.
\xyoption{2cell}
\UseAllTwocells

% We use multicol for the list of chapters between chapters
\usepackage{multicol}

% This is generally recommended for better output
\usepackage{lmodern}
\usepackage[T1]{fontenc}

% For cross-file-references
\usepackage{xr-hyper}

% Package for hypertext links:
\usepackage{hyperref}

% For any local file, say "hello.tex" you want to link to please
% use \externaldocument[hello-]{hello}
\externaldocument[introduction-]{introduction}
\externaldocument[conventions-]{conventions}
\externaldocument[sets-]{sets}
\externaldocument[categories-]{categories}
\externaldocument[topology-]{topology}
\externaldocument[sheaves-]{sheaves}
\externaldocument[sites-]{sites}
\externaldocument[stacks-]{stacks}
\externaldocument[fields-]{fields}
\externaldocument[algebra-]{algebra}
\externaldocument[brauer-]{brauer}
\externaldocument[homology-]{homology}
\externaldocument[derived-]{derived}
\externaldocument[simplicial-]{simplicial}
\externaldocument[more-algebra-]{more-algebra}
\externaldocument[smoothing-]{smoothing}
\externaldocument[modules-]{modules}
\externaldocument[sites-modules-]{sites-modules}
\externaldocument[injectives-]{injectives}
\externaldocument[cohomology-]{cohomology}
\externaldocument[sites-cohomology-]{sites-cohomology}
\externaldocument[dga-]{dga}
\externaldocument[dpa-]{dpa}
\externaldocument[sdga-]{sdga}
\externaldocument[hypercovering-]{hypercovering}
\externaldocument[schemes-]{schemes}
\externaldocument[constructions-]{constructions}
\externaldocument[properties-]{properties}
\externaldocument[morphisms-]{morphisms}
\externaldocument[coherent-]{coherent}
\externaldocument[divisors-]{divisors}
\externaldocument[limits-]{limits}
\externaldocument[varieties-]{varieties}
\externaldocument[topologies-]{topologies}
\externaldocument[descent-]{descent}
\externaldocument[perfect-]{perfect}
\externaldocument[more-morphisms-]{more-morphisms}
\externaldocument[flat-]{flat}
\externaldocument[groupoids-]{groupoids}
\externaldocument[more-groupoids-]{more-groupoids}
\externaldocument[etale-]{etale}
\externaldocument[chow-]{chow}
\externaldocument[intersection-]{intersection}
\externaldocument[pic-]{pic}
\externaldocument[weil-]{weil}
\externaldocument[adequate-]{adequate}
\externaldocument[dualizing-]{dualizing}
\externaldocument[duality-]{duality}
\externaldocument[discriminant-]{discriminant}
\externaldocument[derham-]{derham}
\externaldocument[local-cohomology-]{local-cohomology}
\externaldocument[algebraization-]{algebraization}
\externaldocument[curves-]{curves}
\externaldocument[resolve-]{resolve}
\externaldocument[models-]{models}
\externaldocument[functors-]{functors}
\externaldocument[equiv-]{equiv}
\externaldocument[pione-]{pione}
\externaldocument[etale-cohomology-]{etale-cohomology}
\externaldocument[proetale-]{proetale}
\externaldocument[relative-cycles-]{relative-cycles}
\externaldocument[more-etale-]{more-etale}
\externaldocument[trace-]{trace}
\externaldocument[crystalline-]{crystalline}
\externaldocument[spaces-]{spaces}
\externaldocument[spaces-properties-]{spaces-properties}
\externaldocument[spaces-morphisms-]{spaces-morphisms}
\externaldocument[decent-spaces-]{decent-spaces}
\externaldocument[spaces-cohomology-]{spaces-cohomology}
\externaldocument[spaces-limits-]{spaces-limits}
\externaldocument[spaces-divisors-]{spaces-divisors}
\externaldocument[spaces-over-fields-]{spaces-over-fields}
\externaldocument[spaces-topologies-]{spaces-topologies}
\externaldocument[spaces-descent-]{spaces-descent}
\externaldocument[spaces-perfect-]{spaces-perfect}
\externaldocument[spaces-more-morphisms-]{spaces-more-morphisms}
\externaldocument[spaces-flat-]{spaces-flat}
\externaldocument[spaces-groupoids-]{spaces-groupoids}
\externaldocument[spaces-more-groupoids-]{spaces-more-groupoids}
\externaldocument[bootstrap-]{bootstrap}
\externaldocument[spaces-pushouts-]{spaces-pushouts}
\externaldocument[spaces-chow-]{spaces-chow}
\externaldocument[groupoids-quotients-]{groupoids-quotients}
\externaldocument[spaces-more-cohomology-]{spaces-more-cohomology}
\externaldocument[spaces-simplicial-]{spaces-simplicial}
\externaldocument[spaces-duality-]{spaces-duality}
\externaldocument[formal-spaces-]{formal-spaces}
\externaldocument[restricted-]{restricted}
\externaldocument[spaces-resolve-]{spaces-resolve}
\externaldocument[formal-defos-]{formal-defos}
\externaldocument[defos-]{defos}
\externaldocument[cotangent-]{cotangent}
\externaldocument[examples-defos-]{examples-defos}
\externaldocument[algebraic-]{algebraic}
\externaldocument[examples-stacks-]{examples-stacks}
\externaldocument[stacks-sheaves-]{stacks-sheaves}
\externaldocument[criteria-]{criteria}
\externaldocument[artin-]{artin}
\externaldocument[quot-]{quot}
\externaldocument[stacks-properties-]{stacks-properties}
\externaldocument[stacks-morphisms-]{stacks-morphisms}
\externaldocument[stacks-limits-]{stacks-limits}
\externaldocument[stacks-cohomology-]{stacks-cohomology}
\externaldocument[stacks-perfect-]{stacks-perfect}
\externaldocument[stacks-introduction-]{stacks-introduction}
\externaldocument[stacks-more-morphisms-]{stacks-more-morphisms}
\externaldocument[stacks-geometry-]{stacks-geometry}
\externaldocument[moduli-]{moduli}
\externaldocument[moduli-curves-]{moduli-curves}
\externaldocument[examples-]{examples}
\externaldocument[exercises-]{exercises}
\externaldocument[guide-]{guide}
\externaldocument[desirables-]{desirables}
\externaldocument[coding-]{coding}
\externaldocument[obsolete-]{obsolete}
\externaldocument[fdl-]{fdl}
\externaldocument[index-]{index}

% Theorem environments.
%
\theoremstyle{plain}
\newtheorem{theorem}[subsection]{Theorem}
\newtheorem{proposition}[subsection]{Proposition}
\newtheorem{lemma}[subsection]{Lemma}

\theoremstyle{definition}
\newtheorem{definition}[subsection]{Definition}
\newtheorem{example}[subsection]{Example}
\newtheorem{exercise}[subsection]{Exercise}
\newtheorem{situation}[subsection]{Situation}

\theoremstyle{remark}
\newtheorem{remark}[subsection]{Remark}
\newtheorem{remarks}[subsection]{Remarks}

\numberwithin{equation}{subsection}

% Macros
%
\def\lim{\mathop{\mathrm{lim}}\nolimits}
\def\colim{\mathop{\mathrm{colim}}\nolimits}
\def\Spec{\mathop{\mathrm{Spec}}}
\def\Hom{\mathop{\mathrm{Hom}}\nolimits}
\def\Ext{\mathop{\mathrm{Ext}}\nolimits}
\def\SheafHom{\mathop{\mathcal{H}\!\mathit{om}}\nolimits}
\def\SheafExt{\mathop{\mathcal{E}\!\mathit{xt}}\nolimits}
\def\Sch{\mathit{Sch}}
\def\Mor{\mathop{\mathrm{Mor}}\nolimits}
\def\Ob{\mathop{\mathrm{Ob}}\nolimits}
\def\Sh{\mathop{\mathit{Sh}}\nolimits}
\def\NL{\mathop{N\!L}\nolimits}
\def\CH{\mathop{\mathrm{CH}}\nolimits}
\def\proetale{{pro\text{-}\acute{e}tale}}
\def\etale{{\acute{e}tale}}
\def\QCoh{\mathit{QCoh}}
\def\Ker{\mathop{\mathrm{Ker}}}
\def\Im{\mathop{\mathrm{Im}}}
\def\Coker{\mathop{\mathrm{Coker}}}
\def\Coim{\mathop{\mathrm{Coim}}}

% Boxtimes
%
\DeclareMathSymbol{\boxtimes}{\mathbin}{AMSa}{"02}

%
% Macros for moduli stacks/spaces
%
\def\QCohstack{\mathcal{QC}\!\mathit{oh}}
\def\Cohstack{\mathcal{C}\!\mathit{oh}}
\def\Spacesstack{\mathcal{S}\!\mathit{paces}}
\def\Quotfunctor{\mathrm{Quot}}
\def\Hilbfunctor{\mathrm{Hilb}}
\def\Curvesstack{\mathcal{C}\!\mathit{urves}}
\def\Polarizedstack{\mathcal{P}\!\mathit{olarized}}
\def\Complexesstack{\mathcal{C}\!\mathit{omplexes}}
% \Pic is the operator that assigns to X its picard group, usage \Pic(X)
% \Picardstack_{X/B} denotes the Picard stack of X over B
% \Picardfunctor_{X/B} denotes the Picard functor of X over B
\def\Pic{\mathop{\mathrm{Pic}}\nolimits}
\def\Picardstack{\mathcal{P}\!\mathit{ic}}
\def\Picardfunctor{\mathrm{Pic}}
\def\Deformationcategory{\mathcal{D}\!\mathit{ef}}
\usepackage{cleveref}
\crefname{originallemma}{original lemma}{original lemmas}
\Crefname{originallemma}{Original lemma}{Original lemmas}
\makeatletter
\newlabel{lemma-countable-affine-formal-algebraic-space}{{0AII}{source}{}{sourcepointer-1}{}}
\newlabel{lemma-countable-affine-formal-algebraic-space@cref}{{[originallemma][0][]0AII}{[1][0][]source}{}{}{}}
\makeatother
\crefname{originallemma}{original lemma}{original lemmas}
\Crefname{originallemma}{Original lemma}{Original lemmas}
\makeatletter
\newlabel{lemma-implications-between-types}{{0AIK}{source}{}{sourcepointer-2}{}}
\newlabel{lemma-implications-between-types@cref}{{[originallemma][0][]0AIK}{[1][0][]source}{}{}{}}
\makeatother
\crefname{originallemma}{original lemma}{original lemmas}
\Crefname{originallemma}{Original lemma}{Original lemmas}
\makeatletter
\newlabel{lemma-mcquillan-affine-formal-algebraic-space}{{0AIC}{source}{}{sourcepointer-3}{}}
\newlabel{lemma-mcquillan-affine-formal-algebraic-space@cref}{{[originallemma][0][]0AIC}{[1][0][]source}{}{}{}}
\makeatother
\crefname{originalsection}{original section}{original sections}
\Crefname{originalsection}{Original section}{Original sections}
\makeatletter
\newlabel{more-algebra-section-topological-ring}{{07E7}{source}{}{sourcepointer-4}{}}
\newlabel{more-algebra-section-topological-ring@cref}{{[originalsection][0][]07E7}{[1][0][]source}{}{}{}}
\makeatother
\title{Five equivalent descriptions of a countably indexed affine formal algebraic space}
\author{The Stacks Project Authors}
\date{stacks-project:formal-spaces.tex:local-snapshot-20261012; GNU FDL 1.2 or later}
\begin{document}
\maketitle
\noindent\textit{Editorial note.} All delivered mathematical text is the author's own, copied from the pinned Stacks Project chapter file formal-spaces.tex; no statement or proof step was rewritten, reordered or removed.\par\medskip
\noindent\textit{Editorial note.} The proof uses the author's informal phrase 'etc, etc' in the middle of the nilpotence induction; it lies inside the counted span and is reproduced verbatim, and the induction it abbreviates is recorded as an obligation of the card.\par\medskip
\noindent\textit{Editorial note.} The proof displays one sequence of ideals; it lies inside the counted span and is reproduced verbatim.\par\medskip
\noindent\textit{Editorial note.} The counted spans carry 4 navigation stubs; no bibliography citation occurs in the spans.\par\medskip
\noindent\textit{Editorial note.} The preamble span points at the chapter's own preamble.tex shipped with the frozen source, used verbatim as staged.\par\medskip
\noindent\textit{Allowed standard prerequisites (editorial).}\begin{itemize}
\item The types of affine formal algebraic spaces, in particular the implications between them. Same chapter, Lemma lemma-implications-between-types, invoked at proof lines 2515-2516. Writes the space as the formal spectrum of a weakly admissible topological algebra.
\item Presentations of an affine formal algebraic space as a filtered colimit of affine schemes. Same chapter, Lemma lemma-countable-affine-formal-algebraic-space part (1), invoked at proof lines 2517-2519, and Lemma lemma-mcquillan-affine-formal-algebraic-space, invoked at proof lines 2521-2525. Supply the quotients by weak ideals of definition and the fact that they form a fundamental system of open neighbourhoods of zero.
\item The limit topology on a limit of quotients is linear, complete and its open ideals give a fundamental system. More on Algebra, Section more-algebra-section-topological-ring, invoked at proof lines 2538-2540. Fixes the topological properties used in the nilpotence argument.
\item Weakly admissible topological algebras and weak ideals of definition. Same chapter, the definitions of these notions used throughout at proof lines 2515-2534 and 2544-2549. Fix the properties being proved and the ideals appearing in each description.
\end{itemize}\medskip
\ifdefined\hypertarget\hypertarget{sourcepointer-1}{}\fi
\noindent\textit{Source reference.} Original lemma 0AII: sources/stacks-formal-spaces/pinned-original-source.tex, source lines 2423--2423. Pointer outside the counted span, invoked once for context; the load-bearing argument of the counted statement is the author's proof reproduced in full in the proof span.\par\medskip
\ifdefined\hypertarget\hypertarget{sourcepointer-2}{}\fi
\noindent\textit{Source reference.} Original lemma 0AIK: sources/stacks-formal-spaces/pinned-original-source.tex, source lines 2453--2453. Pointer outside the counted span, invoked once for context; the load-bearing argument of the counted statement is the author's proof reproduced in full in the proof span.\par\medskip
\ifdefined\hypertarget\hypertarget{sourcepointer-3}{}\fi
\noindent\textit{Source reference.} Original lemma 0AIC: sources/stacks-formal-spaces/pinned-original-source.tex, source lines 2135--2135. Pointer outside the counted span, invoked once for context; the load-bearing argument of the counted statement is the author's proof reproduced in full in the proof span.\par\medskip
\ifdefined\hypertarget\hypertarget{sourcepointer-4}{}\fi
\noindent\textit{Source reference.} Original section 07E7: sources/stacks-chow/tags.txt, source lines 9437--9437. Pointer outside the counted span, invoked once for context; the load-bearing argument of the counted statement is the author's proof reproduced in full in the proof span.\par\medskip
\begin{lemma}
\label{lemma-countably-indexed}
Let $S$ be a scheme. Let $X$ be a presheaf on $(\Sch/S)_{fppf}$.
The following are equivalent
\begin{enumerate}
\item $X$ is a countably indexed affine formal algebraic space,
\item $X = \text{Spf}(A)$ where $A$ is a weakly admissible topological
$S$-algebra which has a countable fundamental system of neighbourhoods of $0$,
\item $X = \text{Spf}(A)$ where $A$ is a weakly admissible topological
$S$-algebra which has a fundamental system
$A \supset I_1 \supset I_2 \supset I_3 \supset \ldots$
of weak ideals of definition,
\item $X = \text{Spf}(A)$ where $A$ is a complete topological $S$-algebra
with a fundamental system of open neighbourhoods of $0$ given by a
countable sequence $A \supset I_1 \supset I_2 \supset I_3 \supset \ldots$
of ideals such that $I_n/I_{n + 1}$ is locally nilpotent, and
\item $X = \text{Spf}(A)$ where $A = \lim B/J_n$ with the limit topology
where $B \supset J_1 \supset J_2 \supset J_3 \supset \ldots$ is a
sequence of ideals in an $S$-algebra $B$ with $J_n/J_{n + 1}$
locally nilpotent.
\end{enumerate}
\end{lemma}

\begin{proof}
Assume (1). By Lemma \ref{lemma-implications-between-types}
we can write $X = \text{Spf}(A)$ where $A$ is a weakly admissible
topological $S$-algebra. For any presentation $X = \colim X_n$ as in
Lemma \ref{lemma-countable-affine-formal-algebraic-space} part (1)
we see that $A = \lim A_n$ with $X_n = \Spec(A_n)$ and
$A_n = A/I_n$ for some weak ideal of definition $I_n \subset A$.
This follows from the final statement of
Lemma \ref{lemma-mcquillan-affine-formal-algebraic-space}
which moreover implies that $\{I_n\}$ is a fundamental system
of open neighbourhoods of $0$. Thus we have a sequence
$$
A \supset I_1 \supset I_2 \supset I_3 \supset \ldots
$$
of weak ideals of definition with $A = \lim A/I_n$. In this way
we see that condition (1) implies each of the conditions (2) -- (5).

\medskip\noindent
Assume (5). First note that the limit topology on
$A = \lim B/J_n$ is a linearly topologized, complete topology, see
More on Algebra, Section \ref{more-algebra-section-topological-ring}.
If $f \in A$ maps to zero in $B/J_1$, then some power maps to zero
in $B/J_2$ as its image in $J_1/J_2$ is nilpotent, then a further
power maps to zero in $J_2/J_3$, etc, etc. In this way we see
the open ideal $\Ker(A \to B/J_1)$ is a weak ideal of definition.
Thus $A$ is weakly admissible. In this way we see that (5) implies (2).

\medskip\noindent
It is clear that (4) is a special case of (5) by taking $B = A$.
It is clear that (3) is a special case of (2).

\medskip\noindent
Assume $A$ is as in (2). Let $E_n$ be a countable fundamental
system of neighbourhoods of $0$ in $A$. Since $A$ is a weakly
admissible topological ring we can find open ideals $I_n \subset E_n$.
We can also choose a weak ideal of definition $J \subset A$.
Then $J \cap I_n$ is a fundamental system of weak ideals of definition
of $A$ and we get
$X = \text{Spf}(A) = \colim \Spec(A/(J \cap I_n))$
which shows that $X$ is a countably indexed affine formal algebraic space.
\end{proof}

\end{document}
